% ============================================================
%  §10  Discussion
% ============================================================

\section{Discussion}
\label{sec:discussion}

\subsection{Relation to Existing Frameworks}

\textbf{Dynamical systems approaches to cognition.}
The treatment of concepts as attractors in a dynamical system has
precedent in connectionist and embodied cognition research
\citep[e.g.,][]{thelen1994dynamic}.
The present framework differs in grounding the attractor geometry in
a field-theoretic substrate built from the single foundational axiom of
VED and the effective dynamical conventions of IFGT, and in providing
an explicit account of how discrete linguistic tokens interact with the
continuous field.

\textbf{Distributional semantics.}
Word embedding models place tokens in a geometric space where
proximity reflects co-occurrence patterns.
The conceptual landscape framework is structurally related but
differs fundamentally: in distributional semantics, the geometry
is derived from token co-occurrence statistics; here, it emerges
from the dynamics of token--field coupling.
The geometry reflects the history of \emph{use}, not the statistics
of \emph{co-occurrence}.

\textbf{Predictive processing and free energy.}
Predictive processing accounts treat cognition as the minimization
of prediction error \citep{friston2010free}.
The present framework is complementary: $\sigma_{\mathrm{drive}}$
corresponds to prediction error influx, and the deepening of $\Phi$
wells can be understood as a long-term form of error minimization.
However, the present framework does not commit to a Bayesian
generative model; the field dynamics are more general.

\subsection{Open Questions}

\textbf{Quantitative correspondence.}
The framework establishes structural conditions but does not specify
the concrete form of the convolution kernel $K$ or the parameter
values $D$, $\mu$, $\alpha$, $\Gamma$, $\kappa_B$.
Identifying these from empirical data --- behavioral, linguistic,
or neural --- is a primary open question.

\textbf{Topology and concept individuation.}
The framework defines concepts as locally stable regions of $\Phi$,
but the conditions under which a region counts as a single concept
rather than a cluster of related concepts are not yet specified.
A topological criterion for concept individuation is needed.

\textbf{Multi-agent synchronization dynamics.}
The social landscape is described qualitatively; a formal treatment
of multi-agent token--field interaction, including conditions for
stable shared landscape formation and breakdown, remains to be
developed.

\textbf{Connection to neural implementation.}
The correspondence between the field quantities $I$, $\Phi$, $\JI$
and measurable neural quantities (firing rates, connectivity weights,
oscillatory dynamics) is a major open question that requires
collaboration with empirical neuroscience.
